\chapter{Программный код\label{appendix-extra-examples}}							% 

\begin{verbatim}
	##LoopStatisticsAlgorithm
package algorithms;

import java.util.ArrayList;
import java.util.Collections;
import java.util.List;

public class LoopStatisticsAlgorithm implements StatisticsAlgorithm {

    @Override
    public String getName() {
        return "Loop-based (Imperative)";
    }

    @Override
    public double min(List<Double> data) {
        if (data.isEmpty()) throw new RuntimeException("Пустая выборка");
        double min = Double.MAX_VALUE;
        for (double x : data) {
            if (x < min) min = x;
        }
        return min;
    }

    @Override
    public double max(List<Double> data) {
        if (data.isEmpty()) throw new RuntimeException("Пустая выборка");
        double max = -Double.MAX_VALUE;
        for (double x : data) {
            if (x > max) max = x;
        }
        return max;
    }

    @Override
    public double mean(List<Double> data) {
        if (data.isEmpty()) throw new RuntimeException("Пустая выборка");
        double sum = 0.0;
        for (double x : data) {
            sum += x;
        }
        return sum / data.size();
    }

    @Override
    public double median(List<Double> data) {
        if (data.isEmpty()) throw new RuntimeException("Пустая выборка");
        List<Double> temp = new ArrayList<>(data);
        Collections.sort(temp);
        int n = temp.size();
        int mid = n / 2;
        if (n % 2 == 0) {
            return (temp.get(mid - 1) + temp.get(mid)) / 2.0;
        } else {
            return temp.get(mid);
        }
    }

    @Override
    public double variance(List<Double> data) {
        if (data.isEmpty()) throw new RuntimeException("Пустая выборка");
        if (data.size() < 2) throw new RuntimeException("Недостаточно данных");
        double m = mean(data);
        double sum = 0.0;
        for (double x : data) {
            sum += (x - m) * (x - m);
        }
        return sum / (data.size() - 1);
    }

    @Override
    public double skewness(List<Double> data) {
        if (data.isEmpty()) throw new RuntimeException("Пустая выборка");
        double m = mean(data);
        double v = variance(data);
        double sum = 0.0;
        for (double x : data) {
            sum += Math.pow(x - m, 3);
        }
        return sum / data.size() / Math.pow(v, 1.5);
    }

    @Override
    public double kurtosis(List<Double> data) {
        int n = data.size();
        if (n < 4) throw new RuntimeException("Недостаточно данных для эксцесса");
        double sum1 = 0.0;
        for (double x : data) sum1 += x;
        double mean = sum1 / n;
        double sx2 = 0.0;
        double sum4 = 0.0;
        for (double x : data) {
            double d = x - mean;
            sx2 += d * d;
            sum4 += d * d * d * d;
        }
        sx2 /= (n - 1);
        if (sx2 == 0.0) throw new RuntimeException("Дисперсия равна нулю");
        double sum4Norm = sum4 / Math.pow(sx2, 2);
        return (n * (n + 1.0) / ((n - 1) * (n - 2) * (n - 3))) * sum4Norm
                - 3.0 * (n - 1) * (n - 1) / ((n - 2) * (n - 3));
    }
	

##StatisticsAlgorithm
package algorithms;

import java.util.List;

public interface StatisticsAlgorithm {

    double min(List<Double> data);

    double max(List<Double> data);

    double mean(List<Double> data);

    double median(List<Double> data);

    double variance(List<Double> data);

    double skewness(List<Double> data);

    double kurtosis(List<Double> data);

    String getName();
}

##StreamStatisticsAlgorithm
package algorithms;

import java.util.ArrayList;
import java.util.Collections;
import java.util.List;

public class StreamStatisticsAlgorithm implements StatisticsAlgorithm {

    @Override
    public String getName() {
        return "Stream-based (Functional)";
    }

    @Override
    public double min(List<Double> data) {
        if (data.isEmpty()) throw new RuntimeException("Пустая выборка");
        return data.stream().mapToDouble(Double::doubleValue).min().orElseThrow();
    }

    @Override
    public double max(List<Double> data) {
        if (data.isEmpty()) throw new RuntimeException("Пустая выборка");
        return data.stream().mapToDouble(Double::doubleValue).max().orElseThrow();
    }

    @Override
    public double mean(List<Double> data) {
        if (data.isEmpty()) throw new RuntimeException("Пустая выборка");
        return data.stream().mapToDouble(Double::doubleValue).average().orElseThrow();
    }

    @Override
    public double median(List<Double> data) {
        if (data.isEmpty()) throw new RuntimeException("Пустая выборка");
        List<Double> temp = new ArrayList<>(data);
        Collections.sort(temp);
        int n = temp.size();
        int mid = n / 2;
        if (n % 2 == 0) {
            return (temp.get(mid - 1) + temp.get(mid)) / 2.0;
        } else {
            return temp.get(mid);
        }
    }

    @Override
    public double variance(List<Double> data) {
        if (data.isEmpty()) throw new RuntimeException("Пустая выборка");
        if (data.size() < 2) throw new RuntimeException("Недостаточно данных");
        double m = mean(data);
        double sum = data.stream()
                .mapToDouble(x -> (x - m) * (x - m))
                .sum();
        return sum / (data.size() - 1);
    }

    @Override
    public double skewness(List<Double> data) {
        if (data.isEmpty()) throw new RuntimeException("Пустая выборка");
        double m = mean(data);
        double v = variance(data);
        double sum = data.stream()
                .mapToDouble(x -> Math.pow(x - m, 3))
                .sum();
        return sum / data.size() / Math.pow(v, 1.5);
    }

    @Override
    public double kurtosis(List<Double> data) {
        int n = data.size();
        if (n < 4) throw new RuntimeException("Недостаточно данных для эксцесса");
        double mean = data.stream().mapToDouble(Double::doubleValue).average().orElseThrow();
        double sx2 = data.stream()
                .mapToDouble(x -> (x - mean) * (x - mean))
                .sum() / (n - 1);
        if (sx2 == 0.0) throw new RuntimeException("Дисперсия равна нулю");
        double sum4 = data.stream()
                .mapToDouble(x -> {
                    double d = x - mean;
                    return d * d * d * d;
                })
                .sum();
        double sum4Norm = sum4 / Math.pow(sx2, 2);
        return (n * (n + 1.0) / ((n - 1) * (n - 2) * (n - 3))) * sum4Norm
                - 3.0 * (n - 1) * (n - 1) / ((n - 2) * (n - 3));
    }


##SampleData
package json;

import java.util.List;

public class SampleData {

    public String sampleType;
    public List<Double> data;
    public double minVal;
    public double maxVal;
    public double meanVal;
    public double stddevVal;

    public SampleData() {
    }

    public SampleData(String sampleType, List<Double> data) {
        this.sampleType = sampleType;
        this.data = data;
    }

    public SampleData withUniformParams(double min, double max) {
        this.minVal = min;
        this.maxVal = max;
        return this;
    }

    public SampleData withNormalParams(double mean, double stddev) {
        this.meanVal = mean;
        this.stddevVal = stddev;
        return this;
    }
}

##ConsoleSample
package samples;

import json.SampleData;

import java.util.ArrayList;
import java.util.List;
import java.util.Scanner;

public class ConsoleSample extends Sample {

    public ConsoleSample() {
        super();
    }

    @Override
    public void generate(int n) {
        if (n <= 0) {
            throw new IllegalArgumentException("Размер выборки должен быть положительным");
        }

        data.clear();
        Scanner scanner = new Scanner(System.in);

        System.out.println("Введите " + n + " чисел выборки (каждое с новой строки):");

        int count = 0;
        while (count < n) {
            if (scanner.hasNextDouble()) {
                double value = scanner.nextDouble();
                data.add(value);
                count++;
            } else {
                String invalid = scanner.next();
                System.out.println("Ошибка ввода. Введите число ещё раз:");
            }
        }
    }

    public void setData(List<Double> values) {
        data.clear();
        data.addAll(values);
    }

    @Override
    protected SampleData createSampleData() {
        return new SampleData("ConsoleSample", new ArrayList<>(data));
    }

    @Override
    protected void restoreFromSampleData(SampleData sampleData) {
    }

    @Override
    public String toString() {
        return String.format("ConsoleSample[size=%d]", data.size());
    }

    @Override
    public boolean equals(Object obj) {
        if (this == obj) return true;
        if (obj == null || getClass() != obj.getClass()) return false;
        ConsoleSample that = (ConsoleSample) obj;
        return data.equals(that.data);
    }

    @Override
    public int hashCode() {
        return data.hashCode();
    }
}

##FileSample
package samples;

import json.SampleData;

import java.util.ArrayList;
import java.util.List;

public class FileSample extends Sample {

    public FileSample() {
        super();
    }

    public FileSample(List<Double> values) {
        super();
        this.data.addAll(values);
    }

    @Override
    public void generate(int n) {
    }

    public void setData(List<Double> values) {
        data.clear();
        data.addAll(values);
    }

    @Override
    protected SampleData createSampleData() {
        return new SampleData("FileSample", new ArrayList<>(data));
    }

    @Override
    protected void restoreFromSampleData(SampleData sampleData) {
    }

    @Override
    public String toString() {
        return String.format("FileSample[size=%d]", data.size());
    }

    @Override
    public boolean equals(Object obj) {
        if (this == obj) return true;
        if (obj == null || getClass() != obj.getClass()) return false;
        FileSample that = (FileSample) obj;
        return data.equals(that.data);
    }

    @Override
    public int hashCode() {
        return data.hashCode();
    }
}

##Main
package samples;

import algorithms.LoopStatisticsAlgorithm;
import algorithms.StreamStatisticsAlgorithm;
import json.*;
import utils.*;
import samples.*;

public class Main {

    public static void main(String[] args) {
        UniformSample uniformSample = null;
        NormalSample normalSample = null;
        FileSample fileSample = null;
        ConsoleSample consoleSample = null;

        while (true) {
           System.out.println("\n=====================================");
            System.out.println("       Статистический калькулятор");
            System.out.println("=====================================");
            System.out.println("1. Ввести выборку вручную");
            System.out.println("2. Создать равномерную выборку");
            System.out.println("3. Создать нормальную выборку");
            System.out.println("4. Загрузить из JSON");
            System.out.println("5. Вывести показатели");
            System.out.println("6. Сохранить в JSON");
            System.out.println("7. Оценка взаимозависимости двух выборок");
            System.out.println("8. Выход");
            System.out.println("=====================================");

            int choice = ReadFunctions.readInt("Выберите пункт меню: ", 0, 10);

            switch (choice) {
                case 0: {
                    System.out.println("\n----ТЕСТИРОВАНИЕ----");
                    System.out.println("Запустите тесты через Maven: mvn test");
                    break;
                }
                case 1: {
                    int n = ReadFunctions.readInt("Введите количество элементов: ", 1, 100000);
                    consoleSample = new ConsoleSample();
                    consoleSample.generate(n);
                    System.out.println("\nВыборка сохранена");
                    break;
                }
                case 2: {
                    int n = ReadFunctions.readInt("Размер выборки: ", 1, 100000);
                    double a = ReadFunctions.readDouble("Минимум: ");
                    double b = ReadFunctions.readDouble("Максимум: ", a, Double.MAX_VALUE);
                    uniformSample = new UniformSample(a, b);
                    uniformSample.generate(n);
                    System.out.println("Равномерная выборка создана, n = " + n);
                    break;
                }
                case 3: {
                    int n = ReadFunctions.readInt("Размер выборки: ", 1, 100000);
                    double mean = ReadFunctions.readDouble("Среднее значение: ");
                    double stddev = ReadFunctions.readDouble("Стандартное отклонение: ", 0, Double.MAX_VALUE);
                    normalSample = new NormalSample(mean, stddev);
                    normalSample.generate(n);
                    System.out.println("Нормальная выборка создана, n = " + n);
                    break;
                }
                case 4: {
                    try {
                        String filename = ReadFunctions.readString("Введите имя файла (.json): ");
                        SampleData sampleData = JsonUtils.loadFromJson(filename);

                        switch (sampleData.sampleType) {
                            case "UniformSample":
                                uniformSample = new UniformSample(sampleData.minVal, sampleData.maxVal);
                                uniformSample.loadFromJson(filename);
                                System.out.println("Загружено UniformSample: " + uniformSample.getData().size() + " элементов");
                                break;
                            case "NormalSample":
                                normalSample = new NormalSample(sampleData.meanVal, sampleData.stddevVal);
                                normalSample.loadFromJson(filename);
                                System.out.println("Загружено NormalSample: " + normalSample.getData().size() + " элементов");
                                break;
                            case "FileSample":
                            case "ConsoleSample":
                                fileSample = new FileSample();
                                fileSample.loadFromJson(filename);
                                System.out.println("Загружено: " + fileSample.getData().size() + " элементов");
                                break;
                            default:
                                System.out.println("Неизвестный тип выборки: " + sampleData.sampleType);
                        }
                    } catch (RuntimeException e) {
                        System.out.println("Ошибка: " + e.getMessage());
                    }
                    break;
                }
                case 5: {
                    if (consoleSample != null) {
                        System.out.println("\nВведенная выборка:");
                        consoleSample.showCharact();
                    }
                    if (uniformSample != null) {
                        System.out.println("\nРавномерная выборка:");
                        uniformSample.showCharact();
                    }
                    if (normalSample != null) {
                        System.out.println("\nНормальная выборка:");
                        normalSample.showCharact();
                    }
                    if (fileSample != null) {
                        System.out.println("\nВыборка из файла:");
                        fileSample.showCharact();
                    }
                    break;
                }
                case 6: {
                    int sampleChoice = ReadFunctions.readInt(
                            "\nВыберите выборку для сохранения в JSON: \n1 - Равномерная\n2 - Нормальная\n3 - Из файла\n4 - Введенная\nВведите число: ", 1, 4);

                    Sample selectedSample = pickSample(sampleChoice, uniformSample, normalSample, fileSample, consoleSample);

                    if (selectedSample != null) {
                        try {
                            String filename = ReadFunctions.readString("Введите имя файла (.json): ");
                            selectedSample.saveToJson(filename.endsWith(".json") ? filename : filename + ".json");
                            System.out.println("Выборка сохранена в JSON");
                        } catch (RuntimeException e) {
                            System.out.println("Ошибка: " + e.getMessage());
                        }
                    } else {
                        System.out.println("Сначала создайте выборку");
                    }
                    break;
                }
                case 7: {
                    int a = ReadFunctions.readInt(
                            "Выберите первую выборку:\n1 - Равномерная\n2 - Нормальная\n3 - Из файла\n4 - Введенная\nВыбор: ", 1, 4);
                    int b = ReadFunctions.readInt("Выберите вторую выборку: ", 1, 4);

                    Sample s1 = pickSample(a, uniformSample, normalSample, fileSample, consoleSample);
                    Sample s2 = pickSample(b, uniformSample, normalSample, fileSample, consoleSample);

                    if (s1 == null || s2 == null) {
                        System.out.println("Одна из выборок не существует");
                        break;
                    }
                    try {
                        double r = StatsUtils.correlation(s1, s2);
                        System.out.printf("Коэффициент корреляции Пирсона: %.4f%n", r);
                    } catch (RuntimeException e) {
                        System.out.println("Ошибка: " + e.getMessage());
                    }
                    break;
                }

                case 8: {
                    System.out.println("Выход...");
                    return;
                }
                default:
                    System.out.println("Неверный выбор!");
                    break;
            }
        }
    }

    private static Sample pickSample(int choice, UniformSample uniform, NormalSample normal, FileSample file, ConsoleSample console) {
        return switch (choice) {
            case 1 -> uniform;
            case 2 -> normal;
            case 3 -> file;
            case 4 -> console;
            default -> null;
        };
    }
}

##NormalSample
package samples;

import json.SampleData;

import java.util.ArrayList;
import java.util.List;
import java.util.Random;

public class NormalSample extends Sample {

    private double meanVal;
    private double stddevVal;

    private static final Random random = new Random();

    public NormalSample() {
        super();
        this.meanVal = 0.0;
        this.stddevVal = 1.0;
    }

    public NormalSample(double mean, double stddev) {
        super();
        this.meanVal = mean;
        this.stddevVal = stddev;
    }

    @Override
    public void generate(int n) {
        if (n <= 0) {
            throw new IllegalArgumentException("Размер выборки должен быть положительным");
        }
        if (stddevVal <= 0) {
            throw new IllegalArgumentException("Стандартное отклонение должно быть положительным");
        }

        data.clear();

        if (data instanceof ArrayList) {
            ((ArrayList<Double>) data).ensureCapacity(n);
        }

        for (int i = 0; i < n; i++) {
            double randomValue = meanVal + stddevVal * random.nextGaussian();
            data.add(randomValue);
        }
    }

    @Override
    protected SampleData createSampleData() {
        return new SampleData("NormalSample", new ArrayList<>(data))
                .withNormalParams(meanVal, stddevVal);
    }

    @Override
    protected void restoreFromSampleData(SampleData sampleData) {
        this.meanVal = sampleData.meanVal;
        this.stddevVal = sampleData.stddevVal;
    }

    public double getMeanVal() {
        return meanVal;
    }

    public void setMeanVal(double meanVal) {
        this.meanVal = meanVal;
    }

    public double getStddevVal() {
        return stddevVal;
    }

    public void setStddevVal(double stddevVal) {
        this.stddevVal = stddevVal;
    }

    @Override
    public String toString() {
        return String.format("NormalSample[mean=%.2f, stddev=%.2f, size=%d]",
                meanVal, stddevVal, data.size());
    }

    @Override
    public boolean equals(Object obj) {
        if (this == obj) return true;
        if (obj == null || getClass() != obj.getClass()) return false;
        NormalSample that = (NormalSample) obj;
        return Double.compare(that.meanVal, meanVal) == 0 &&
                Double.compare(that.stddevVal, stddevVal) == 0 &&
                data.equals(that.data);
    }

    @Override
    public int hashCode() {
        int result;
        long temp;
        temp = Double.doubleToLongBits(meanVal);
        result = (int) (temp ^ (temp >>> 32));
        temp = Double.doubleToLongBits(stddevVal);
        result = 31 * result + (int) (temp ^ (temp >>> 32));
        result = 31 * result + data.hashCode();
        return result;
    }
}

##Sample
package samples;

import algorithms.StatisticsAlgorithm;
import algorithms.LoopStatisticsAlgorithm;
import json.SampleData;
import utils.JsonUtils;

import java.io.*;
import java.nio.file.*;
import java.util.*;

public abstract class Sample {

    protected List<Double> data;
    private StatisticsAlgorithm algorithm = new LoopStatisticsAlgorithm();

    public Sample() {
        this.data = new ArrayList<>();
    }

    public Sample(StatisticsAlgorithm algorithm) {
        this.data = new ArrayList<>();
        this.algorithm = algorithm;
    }

    public void setData(List<Double> values){
        this.data.clear();
        this.data.addAll(values);
    }

    public abstract void generate(int n);

    public List<Double> getData() {
        return Collections.unmodifiableList(data);
    }

    public void setAlgorithm(StatisticsAlgorithm algorithm) {
        this.algorithm = algorithm;
    }

    public StatisticsAlgorithm getAlgorithm() {
        return this.algorithm;
    }

    public double min() {
        return algorithm.min(data);
    }

    public double max() {
        return algorithm.max(data);
    }

    public double mean() {
        return algorithm.mean(data);
    }

    public double median() {
        return algorithm.median(data);
    }

    public double variance() {
        return algorithm.variance(data);
    }

    public double skewness() {
        return algorithm.skewness(data);
    }

    public double kurtosis() {
        return algorithm.kurtosis(data);
    }

    public List<Double> quartiles() {
        if (data.isEmpty()) {
            throw new RuntimeException("Пустая выборка");
        }

        List<Double> temp = new ArrayList<>(data);
        Collections.sort(temp);

        int n = temp.size();
        double q1 = getQuantile(temp, 0.25);
        double q2 = median();
        double q3 = getQuantile(temp, 0.75);

        return Arrays.asList(q1, q2, q3);
    }

    private double getQuantile(List<Double> sortedData, double q) {
        int n = sortedData.size();
        double pos = q * (n - 1);
        int index = (int) pos;
        double frac = pos - index;

        if (index + 1 < n) {
            return sortedData.get(index) + frac * (sortedData.get(index + 1) - sortedData.get(index));
        } else {
            return sortedData.get(index);
        }
    }

    public boolean isNormal(double alpha) {
        double jb = jarqueBera();
        double critical;

        if (alpha == 0.10) critical = 4.605;
        else if (alpha == 0.05) critical = 5.991;
        else if (alpha == 0.01) critical = 9.210;
        else throw new RuntimeException("Неподдерживаемый уровень значимости: " + alpha);

        return jb < critical;
    }

    public double jarqueBera() {
        if (data.size() < 3) {
            throw new RuntimeException("Недостаточно данных для теста Jarque-Bera");
        }
        double n = (double) data.size();
        double s = skewness();
        double k = kurtosis();
        return (n / 6.0) * (s * s + (k * k) / 4.0);
    }

    public void saveToFile(String filename) {
        try (BufferedWriter writer = Files.newBufferedWriter(Paths.get(filename))) {
            for (double x : data) {
                writer.write(String.valueOf(x));
                writer.newLine();
            }
        } catch (IOException e) {
            throw new RuntimeException("Не удалось сохранить файл: " + e.getMessage());
        }
    }

    public void loadFromFile(String filename) {
        data.clear();
        try (BufferedReader reader = Files.newBufferedReader(Paths.get(filename))) {
            String line;
            while ((line = reader.readLine()) != null) {
                if (!line.trim().isEmpty()) {
                    data.add(Double.parseDouble(line.trim()));
                }
            }
        } catch (IOException e) {
            throw new RuntimeException("Не удалось загрузить файл: " + e.getMessage());
        }

        if (data.isEmpty()) {
            throw new RuntimeException("Файл пуст или не содержит чисел");
        }
    }

    public void saveToJson(String filename) {
        SampleData sampleData = createSampleData();
        JsonUtils.saveToJson(sampleData, filename);
    }

    public void loadFromJson(String filename) {
        SampleData sampleData = JsonUtils.loadFromJson(filename);
        data.clear();
        data.addAll(sampleData.data);
        restoreFromSampleData(sampleData);
    }

    protected SampleData createSampleData() {
        return new SampleData(this.getClass().getSimpleName(), new ArrayList<>(data));
    }

    protected void restoreFromSampleData(SampleData sampleData) {
    }

    public void showCharact() {
        System.out.println("=======================================");
        System.out.println("  Статистика выборки: " + this.getClass().getSimpleName());
        System.out.println("  Алгоритм: " + algorithm.getName());
        System.out.println("=======================================");
        System.out.println("  Размер выборки: " + data.size());
        System.out.println("---------------------------------------");
        System.out.printf("  Минимум:        %.4f%n", this.min());
        System.out.printf("  Максимум:       %.4f%n", this.max());
        System.out.printf("  Среднее:        %.4f%n", this.mean());
        System.out.printf("  Медиана:        %.4f%n", this.median());
        System.out.println("---------------------------------------");

        List<Double> q = this.quartiles();
        System.out.printf("  Квартиль 0.25:  %.4f%n", q.get(0));
        System.out.printf("  Квартиль 0.50:  %.4f%n", q.get(1));
        System.out.printf("  Квартиль 0.75:  %.4f%n", q.get(2));
        System.out.println("---------------------------------------");

        System.out.printf("  Дисперсия:      %.4f%n", this.variance());
        System.out.printf("  Асимметрия:     %.4f%n", this.skewness());
        System.out.printf("  Эксцесс:        %.4f%n", this.kurtosis());
        System.out.println("("=======================================");

        try {
            double jb = this.jarqueBera();
            boolean isNormal = this.isNormal(0.05);
            System.out.printf("  Тест Jarque-Bera: %.4f%n", jb);
            System.out.println("  Нормальность: " + (isNormal ? "не отвергается" : "отвергается"));
        } catch (Exception e) {
            System.out.println("  Тест на нормальность: невозможно проверить");
        }
        System.out.println("("=======================================");
    }

    @Override
    public String toString() {
        return String.format("%s[size=%d, algorithm=%s]",
                this.getClass().getSimpleName(),
                data.size(),
                algorithm.getName());
    }

    @Override
    public boolean equals(Object obj) {
        if (this == obj) return true;
        if (obj == null || getClass() != obj.getClass()) return false;
        Sample sample = (Sample) obj;
        return data.equals(sample.data);
    }

    @Override
    public int hashCode() {
        return data.hashCode();
    }
}

##UniformSample
package samples;

import json.SampleData;

import java.util.ArrayList;
import java.util.List;

public class UniformSample extends Sample {

    private double minVal;
    private double maxVal;

    public UniformSample() {
        super();
        this.minVal = 0.0;
        this.maxVal = 1.0;
    }

    public UniformSample(double min, double max) {
        super();
        this.minVal = min;
        this.maxVal = max;
    }

    @Override
    public void generate(int n) {
        if (n <= 0) {
            throw new IllegalArgumentException("Размер выборки должен быть положительным");
        }
        if (minVal >= maxVal) {
            throw new IllegalArgumentException("minVal должен быть меньше maxVal");
        }

        data.clear();

        if (data instanceof ArrayList) {
            ((ArrayList<Double>) data).ensureCapacity(n);
        }

        for (int i = 0; i < n; i++) {
            double randomValue = minVal + (maxVal - minVal) * Math.random();
            data.add(randomValue);
        }
    }

    @Override
    protected SampleData createSampleData() {
        return new SampleData("UniformSample", new ArrayList<>(data))
                .withUniformParams(minVal, maxVal);
    }

    @Override
    protected void restoreFromSampleData(SampleData sampleData) {
        this.minVal = sampleData.minVal;
        this.maxVal = sampleData.maxVal;
    }

    public double getMinVal() {
        return minVal;
    }

    public void setMinVal(double minVal) {
        this.minVal = minVal;
    }

    public double getMaxVal() {
        return maxVal;
    }

    public void setMaxVal(double maxVal) {
        this.maxVal = maxVal;
    }

    @Override
    public String toString() {
        return String.format("UniformSample[min=%.2f, max=%.2f, size=%d]",
                minVal, maxVal, data.size());
    }

    @Override
    public boolean equals(Object obj) {
        if (this == obj) return true;
        if (obj == null || getClass() != obj.getClass()) return false;
        UniformSample that = (UniformSample) obj;
        return Double.compare(that.minVal, minVal) == 0 &&
                Double.compare(that.maxVal, maxVal) == 0 &&
                data.equals(that.data);
    }

    @Override
    public int hashCode() {
        int result;
        long temp;
        temp = Double.doubleToLongBits(minVal);
        result = (int) (temp ^ (temp >>> 32));
        temp = Double.doubleToLongBits(maxVal);
        result = 31 * result + (int) (temp ^ (temp >>> 32));
        result = 31 * result + data.hashCode();
        return result;
    }
}

##JsonUtils
package utils;

import com.google.gson.Gson;
import com.google.gson.GsonBuilder;
import com.google.gson.stream.JsonReader;
import json.SampleData;

import java.io.*;
import java.nio.file.*;

public class JsonUtils {

    private static final Gson gson = new GsonBuilder()
            .setPrettyPrinting()
            .create();

    private JsonUtils() {
    }

    public static void saveToJson(SampleData sampleData, String filename) {
        try (BufferedWriter writer = Files.newBufferedWriter(Paths.get(filename))) {
            gson.toJson(sampleData, writer);
        } catch (IOException e) {
            throw new RuntimeException("Не удалось сохранить JSON: " + e.getMessage());
        }
    }

    public static SampleData loadFromJson(String filename) {
        try (BufferedReader reader = Files.newBufferedReader(Paths.get(filename));
             JsonReader jsonReader = new JsonReader(reader)) {
            return gson.fromJson(jsonReader, SampleData.class);
        } catch (IOException e) {
            throw new RuntimeException("Не удалось загрузить JSON: " + e.getMessage());
        }
    }

    public static String toJson(SampleData sampleData) {
        return gson.toJson(sampleData);
    }
}

##ReadFunctions
package utils;

import java.util.Scanner;

public class ReadFunctions {

    private static Scanner scanner = new Scanner(System.in);


    public static Scanner getScanner(){
        return scanner;
    }

    private ReadFunctions() {
    }

    public static int readInt(String prompt, int min, int max) {
        int value;
        while (true) {
            System.out.print(prompt);
            if (scanner.hasNextInt()) {
                value = scanner.nextInt();
                if (value >= min && value <= max) {
                    return value;
                } else {
                    System.out.println("Число должно быть от " + min + " до " + max);
                }
            } else {
                System.out.println("Неверный ввод. Попробуйте еще раз");
                scanner.next();
            }
        }
    }

    public static int readInt(String prompt) {
        return readInt(prompt, Integer.MIN_VALUE, Integer.MAX_VALUE);
    }

    public static double readDouble(String prompt, double min, double max) {
        double value;
        while (true) {
            System.out.print(prompt);
            if (scanner.hasNextDouble()) {
                value = scanner.nextDouble();
                if (value >= min && value <= max) {
                    return value;
                } else {
                    System.out.println("Число должно быть от " + min + " до " + max);
                }
            } else {
                System.out.println("Неверный ввод. Попробуйте еще раз");
                scanner.next();
            }
        }
    }

    public static double readDouble(String prompt) {
        return readDouble(prompt, Double.NEGATIVE_INFINITY, Double.POSITIVE_INFINITY);
    }

    public static String readString(String prompt) {
        System.out.print(prompt);
        return scanner.next();
    }
}

##StatsUtils
package utils;

import samples.Sample;

import java.util.List;

public class StatsUtils {

    private StatsUtils() {
    }

    public static double correlation(Sample a, Sample b) {
        List<Double> x = a.getData();
        List<Double> y = b.getData();

        if (x.size() != y.size()) {
            throw new RuntimeException("Выборки должны быть одного размера");
        }
        if (x.size() < 2) {
            throw new RuntimeException("Недостаточно данных");
        }

        double meanX = a.mean();
        double meanY = b.mean();

        double num = 0.0;
        double sx = 0.0;
        double sy = 0.0;

        for (int i = 0; i < x.size(); i++) {
            double dx = x.get(i) - meanX;
            double dy = y.get(i) - meanY;
            num += dx * dy;
            sx += dx * dx;
            sy += dy * dy;
        }

        if (sx == 0.0 || sy == 0.0) {
            throw new RuntimeException("Нулевая дисперсия одной из выборок");
        }

        return num / Math.sqrt(sx * sy);
    }
}

##LoopStatisticsAlgorithmTest
package algorithm;

import algorithms.LoopStatisticsAlgorithm;
import algorithms.StatisticsAlgorithm;
import algorithms.StreamStatisticsAlgorithm;
import org.junit.jupiter.api.Test;
import org.junit.jupiter.api.DisplayName;
import org.junit.jupiter.params.ParameterizedTest;
import org.junit.jupiter.params.provider.ValueSource;
import org.mockito.Mockito;

import static org.junit.jupiter.api.Assertions.*;
import static org.junit.jupiter.api.Assumptions.*;
import static org.hamcrest.MatcherAssert.assertThat;
import static org.hamcrest.Matchers.*;
import static org.mockito.Mockito.*;

import java.util.Arrays;
import java.util.Collections;
import java.util.List;
import java.util.ArrayList;

import static org.mockito.ArgumentMatchers.anyList;


@DisplayName("Тестирование LoopStatisticsAlgorithm")
public class LoopStatisticsAlgorithmTest {
    private final StatisticsAlgorithm loopAlgorithm = new LoopStatisticsAlgorithm();
    private final StatisticsAlgorithm streamAlgorithm = new StreamStatisticsAlgorithm();

    //ГРАНИЧНЫЕ ЗНАЧЕНИЯ (BoundaryValue)
    @Test
    @DisplayName("Граничное значение: пустой список должен выбрасывать исключение")
    void testMinEmptyList_BoundaryValue() {
        //граничное значение 0 элементов
        assertThrows(RuntimeException.class, () -> loopAlgorithm.min(Collections.emptyList()));
    }

    @Test
    @DisplayName("Граничное значение: один элемент")
    void testMinSingleElement_BoundaryValue() {
        //граничное значение 1 элемент
        List<Double> data = Collections.singletonList(5.0);
        double result = loopAlgorithm.min(data);
        assertEquals(5.0, result, 0.0001);
    }

    @Test
    @DisplayName("Граничное значение: 1 элемент для дисперсии должен выбрасывать исключение")
    void testVarianceOneElement_BoundaryValue() {
        //меньше минимального: 1 элемент
        List<Double> data = Collections.singletonList(5.0);
        assertThrows(RuntimeException.class, () -> loopAlgorithm.variance(data));
    }

    @Test
    @DisplayName("Граничное значение: 3 элемента для эксцесса должен выбрасывать исключение")
    void testKurtosisThreeElements_BoundaryValue() {
        //меньше минимального: 3 элемента
        List<Double> data = Arrays.asList(1.0, 2.0, 3.0);
        assertThrows(RuntimeException.class, () -> loopAlgorithm.kurtosis(data));
    }

    // КЛАССЫ ЭКВИВАЛЕНТНОСТИ (Equivalence Partitioning)
    @ParameterizedTest
    @ValueSource(doubles = {1.0, 5.5, 100.0})
    @DisplayName("Класс эквивалентности: положительные числа")
    void testMeanPositiveNumbers_EquivalenceClass(double value) {
        List<Double> data = Arrays.asList(value, value, value);
        double result = loopAlgorithm.mean(data);
        assertThat(result, is(closeTo(value, 0.0001)));
    }

    @ParameterizedTest
    @ValueSource(doubles = {-1.0, -5.5, -100.0})
    @DisplayName("Класс эквивалентности: отрицательные числа")
    void testMeanNegativeNumbers_EquivalenceClass(double value) {
        List<Double> data = Arrays.asList(value, value, value);
        double result = loopAlgorithm.mean(data);
        assertThat(result, is(closeTo(value, 0.0001)));
    }

    @Test
    @DisplayName("Медиана: нечетное количество (3 элемента)")
    void testMedianOddCount() {
        List<Double> data = Arrays.asList(1.0, 2.0, 3.0);
        assertEquals(2.0, loopAlgorithm.median(data), 0.0001);
    }

    @Test
    @DisplayName("Медиана: четное количество (4 элемента)")
    void testMedianEvenCount() {
        List<Double> data = Arrays.asList(1.0, 2.0, 3.0, 4.0);
        assertEquals(2.5, loopAlgorithm.median(data), 0.0001);
    }

    //ПОКРЫТИЕ ОПЕРАТОРОВ (Statement Testing)
    @Test
    @DisplayName("Покрытие операторов: все методы должны выполниться")
    void testAllMethodExecute_StatementCoverage() {
        List<Double> data = Arrays.asList(1.0, 2.0, 3.0, 4.0, 5.0);

        //каждый метод выполнится хотя бы один раз
        double min = loopAlgorithm.min(data);
        double max = loopAlgorithm.max(data);
        double mean = loopAlgorithm.mean(data);
        double median = loopAlgorithm.median(data);
        double variance = loopAlgorithm.variance(data);
        double skewness = loopAlgorithm.skewness(data);
        double kurtosis = loopAlgorithm.kurtosis(data);
        String name = loopAlgorithm.getName();

        //проверка, что все значения получены
        assertAll("Все методы выполнены успешно",
                () -> assertNotNull(min),
                () -> assertNotNull(max),
                () -> assertNotNull(mean),
                () -> assertNotNull(median),
                () -> assertNotNull(variance),
                () -> assertNotNull(skewness),
                () -> assertNotNull(kurtosis),
                () -> assertNotNull(name));
    }

    //ПОКРЫТИЕ ВЕТВЛЕНИЙ (Branch Testing)
    @Test
    @DisplayName("Покрытие ветвлений: медиана для нечетного количества (else)")
    void testMedianOddCount_BranchCoverage() {
        //n%2!=0
        List<Double> data = Arrays.asList(1.0, 3.0, 2.0);
        double result = loopAlgorithm.median(data);
        assertEquals(2.0, result, 0.0001);
    }

    @Test
    @DisplayName("Покрытие ветвлений: медиана для нечетного количества (if)")
    void testMedianEvenCount_BranchCoverage() {
        //n%2==0
        List<Double> data = Arrays.asList(1.0, 3.0, 2.0, 4.0);
        double result = loopAlgorithm.median(data);
        assertEquals(2.5, result, 0.0001);
    }

    @Test
    @DisplayName("Покрытие ветвлений: isEmpty() - true/false")
    void testIsEmpty_BranchCoverage() {
        //пустой список
        assertThrows(RuntimeException.class, () -> loopAlgorithm.min(Collections.emptyList()));

        //заполненный список
        List<Double> data = Collections.singletonList(1.0);
        assertDoesNotThrow(() -> loopAlgorithm.min(data));
    }

    //СРАВНЕНИЕ АЛГОРИТМОВ
    @Test
    @DisplayName("Сравнение алгоритмов: Loop и Stream должны давать одинаковые результаты")
    void testLoopVsStreamAlgorithmComparison() {
        List<Double> data = Arrays.asList(1.0, 5.0, 3.0, 9.0, 2.0, 7.0, 4.0);

        double loopMin = loopAlgorithm.min(data);
        double streamMin = streamAlgorithm.min(data);
        assertEquals(loopMin, streamMin, 0.0001, "min должен совпадать");

        double loopMax = loopAlgorithm.max(data);
        double streamMax = streamAlgorithm.max(data);
        assertEquals(loopMax, streamMax, 0.0001, "max должен совпадать");

        double loopMean = loopAlgorithm.mean(data);
        double streamMean = streamAlgorithm.mean(data);
        assertEquals(loopMean, streamMean, 0.0001, "mean должен совпадать");

        double loopMedian = loopAlgorithm.median(data);
        double streamMedian = streamAlgorithm.median(data);
        assertEquals(loopMedian, streamMedian, 0.0001, "median должен совпадать");

        double loopVariance = loopAlgorithm.variance(data);
        double streamVariance = streamAlgorithm.variance(data);
        assertEquals(loopVariance, streamVariance, 0.0001, "variance должен совпадать");
    }

    @ParameterizedTest
    @ValueSource(ints = {5, 10, 20, 50})
    @DisplayName("Сравнение алгоритмов на разных размерах выборки")
    void testLoopVsStreamVariousSizes(int size) {
        List<Double> data = new ArrayList<>();
        for (int i = 0; i < size; i++) {
            data.add((double) i);
        }

        double loopResult = loopAlgorithm.mean(data);
        double streamResult = streamAlgorithm.mean(data);

        assertEquals(loopResult, streamResult, 0.0001,
                "Результаты должны совпадать для размера" + size);
    }

    //ASSUMPTIONS
    @Test
    @DisplayName("Assumption: тест имеет смысл только для положительных размеров")
    void testWithAssumptionOnSize() {
        List<Double> data = Arrays.asList(1.0, 2.0, 3.0);

        //предположение: тест имеет смысл только если данных больше 2
        assumeTrue(data.size() > 2, "Требуется минимум 3 элемента");

        double variance = loopAlgorithm.variance(data);
        assertTrue(variance >= 0, "Дисперсия не может быть отрицательной");
    }

    @Test
    @DisplayName("Assumption: пропускаем тест если выборка пустая")
    void testWithAssumptionOnEmptyList() {
        List<Double> data = Arrays.asList(1.0, 2.0, 3.0, 4.0, 5.0);

        //предположение: данные не должны быть пустыми
        assumeFalse(data.isEmpty(), "Данные не должны быть пустыми");

        double mean = loopAlgorithm.mean(data);
        assertThat(mean, is(greaterThan(0.0)));
    }

    //MOCKS
    @Test
    @DisplayName("Mock: создание мока StatisticsAlgorithm")
    void testMockStatisticsAlgorithm() {
        //создаем пустой мок
        StatisticsAlgorithm mockAlg = Mockito.mock(StatisticsAlgorithm.class);

        //проверка, что мок создан
        assertNotNull(mockAlg);

        //вызов метода (возврат default для double-0.0)
        double result = mockAlg.mean(Collections.emptyList());
        assertEquals(0.0, result, 0.0001);

    }

    @Test
    @DisplayName("Stub: настройка поведения мок объекта")
    void testStubStatisticsAlgorithm() {
        //stub: мок с настроенным поведением
        StatisticsAlgorithm stubAlg = Mockito.mock(StatisticsAlgorithm.class);

        //настройка возвращаемого значения
        when(stubAlg.mean(anyList())).thenReturn(42.0);
        when(stubAlg.min(anyList())).thenReturn(1.0);
        when(stubAlg.max(anyList())).thenReturn(100.0);

        //проверка возвращаемых значений
        assertEquals(42.0, stubAlg.mean(Collections.emptyList()), 0.0001);
        assertEquals(1.0, stubAlg.min(Collections.emptyList()), 0.0001);
        assertEquals(100.0, stubAlg.max(Collections.emptyList()), 0.0001);

        verify(stubAlg, times(1)).mean(anyList());
        verify(stubAlg, times(1)).min(anyList());
        verify(stubAlg, times(1)).max(anyList());
    }

    @Test
    @DisplayName("Spy: шпион на реальном объекте LoopStatisticsAlgorithm")
    void testSpyStatisticsAlgorithm() {
        //создаем шпиона на реальном объекте
        LoopStatisticsAlgorithm spyAlg = Mockito.spy(new LoopStatisticsAlgorithm());

        List<Double> data = Arrays.asList(1.0, 2.0, 3.0, 4.0, 5.0);

        double mean = spyAlg.mean(data);
        verify(spyAlg, atLeastOnce()).mean(data);
        assertEquals(3.0, mean, 0.0001);
    }

    //MATCHERS
    @Test
    @DisplayName("Matchers: использование Hamcrest матчеров")
    void testWithHamcrestMatchers() {
        List<Double> data = Arrays.asList(1.0, 2.0, 3.0, 4.0, 5.0);

        double mean = loopAlgorithm.mean(data);
        double min = loopAlgorithm.min(data);
        double max = loopAlgorithm.max(data);

        //matcher 1 closeTo
        assertThat(mean, is(closeTo(3.0, 0.0001)));

        //matcher 2 greater/lessthan
        assertThat(min, lessThan(mean));
        assertThat(max, greaterThan(mean));

        //matcher 3 equalto
        assertThat(min, is(equalTo(1.0)));
        assertThat(max, is(equalTo(5.0)));
    }

    //НЕГАТИВНЫЕ СЦЕНАРИИ
    @Test
    @DisplayName("Негативный сценарий: null должен выбрасывать исключение")
    void testWithNullData() {
        assertThrows(NullPointerException.class, () -> loopAlgorithm.mean(null));
    }

    @Test
    @DisplayName("Негативный сценарий: эксцесс при нулевой дисперсии")
    void testKurtosisZeroVariance() {
        List<Double> data = Arrays.asList(5.0, 5.0, 5.0, 5.0); //дисперсия 0, т.к. элементы одинаковые
        assertThrows(RuntimeException.class, () -> loopAlgorithm.kurtosis(data));
    }

    // ==================== МУТАЦИОННЫЕ ТЕСТЫ ====================

    @Test
    @DisplayName("Мутация: getName должен вернуть точную строку")
    void testGetNameExact() {
        assertEquals("Loop-based (Imperative)", loopAlgorithm.getName());
        assertNotEquals("", loopAlgorithm.getName());
        assertNotEquals("Stream-based", loopAlgorithm.getName());
    }

    @Test
    @DisplayName("Мутация: min с дубликатами минимума")
    void testMinWithDuplicateValues() {
        List<Double> data = Arrays.asList(3.0, 1.0, 1.0, 5.0);
        assertEquals(1.0, loopAlgorithm.min(data), 0.0001);
    }

    @Test
    @DisplayName("Мутация: max с дубликатами максимума")
    void testMaxWithDuplicateValues() {
        List<Double> data = Arrays.asList(1.0, 5.0, 5.0, 3.0);
        assertEquals(5.0, loopAlgorithm.max(data), 0.0001);
    }

    @Test
    @DisplayName("Мутация: min граничное значение (0 и отрицательное)")
    void testMinWithZeroAndNegative() {
        List<Double> data = Arrays.asList(0.0, -5.0, 3.0);
        assertEquals(-5.0, loopAlgorithm.min(data), 0.0001);
    }

    @Test
    @DisplayName("Мутация: max граничное значение (0 и отрицательное)")
    void testMaxWithZeroAndNegative() {
        List<Double> data = Arrays.asList(0.0, -5.0, 3.0);
        assertEquals(3.0, loopAlgorithm.max(data), 0.0001);
    }

    @Test
    @DisplayName("Мутация: variance точное значение")
    void testVariancePreciseValue() {
        List<Double> data = Arrays.asList(1.0, 2.0, 3.0, 4.0, 5.0);
        double variance = loopAlgorithm.variance(data);
        assertEquals(2.5, variance, 0.0001);
    }

    @Test
    @DisplayName("Мутация: variance два элемента")
    void testVarianceTwoElements() {
        List<Double> data = Arrays.asList(1.0, 3.0);
        double variance = loopAlgorithm.variance(data);
        assertEquals(2.0, variance, 0.0001);
    }

    @Test
    @DisplayName("Мутация: skewness не пустой, не пусто")
    void testSkewnessNotEmpty() {
        List<Double> data = Arrays.asList(1.0, 2.0, 3.0, 4.0, 5.0);
        double skewness = loopAlgorithm.skewness(data);
        assertEquals(0.0, skewness, 0.0001);  // симметричное распределение
    }

    @Test
    @DisplayName("Мутация: skewness отрицательная асимметрия")
    void testSkewnessNegativeSkew() {
        List<Double> data = Arrays.asList(1.0, 2.0, 3.0, 4.0, 10.0);
        double skewness = loopAlgorithm.skewness(data);
        assertTrue(skewness > 0);  // положительная асимметрия (хвост вправо)
    }

    @Test
    @DisplayName("Мутация: kurtosis граничное значение (4 элемента)")
    void testKurtosisMinElements() {
        List<Double> data = Arrays.asList(1.0, 2.0, 3.0, 4.0);
        double kurtosis = loopAlgorithm.kurtosis(data);
        assertTrue(Double.isFinite(kurtosis));
    }

    @Test
    @DisplayName("Мутация: kurtosis больше элементов")
    void testKurtosisManyElements() {
        List<Double> data = Arrays.asList(1.0, 2.0, 3.0, 4.0, 5.0, 6.0, 7.0, 8.0, 9.0, 10.0);
        double kurtosis = loopAlgorithm.kurtosis(data);
        assertTrue(Double.isFinite(kurtosis));
        assertTrue(Math.abs(kurtosis) < 10);  // должен быть в разумном диапазоне
    }

    @Test
    @DisplayName("Мутация: min отрицательные числа")
    void testMinNegativeNumbers() {
        List<Double> data = Arrays.asList(-5.0, -1.0, -10.0);
        assertEquals(-10.0, loopAlgorithm.min(data), 0.0001);
    }

    @Test
    @DisplayName("Мутация: max отрицательные числа")
    void testMaxNegativeNumbers() {
        List<Double> data = Arrays.asList(-5.0, -1.0, -10.0);
        assertEquals(-1.0, loopAlgorithm.max(data), 0.0001);
    }

    @Test
    @DisplayName("Мутация: median с дублирующимся значением")
    void testMedianWithDuplicates() {
        List<Double> data = Arrays.asList(1.0, 2.0, 2.0, 3.0, 4.0);
        assertEquals(2.0, loopAlgorithm.median(data), 0.0001);
    }

    @Test
    @DisplayName("Мутация: mean точное значение")
    void testMeanPreciseValue() {
        List<Double> data = Arrays.asList(2.0, 4.0, 6.0);
        assertEquals(4.0, loopAlgorithm.mean(data), 0.0001);
    }

    @Test
    @DisplayName("Мутация: variance < 2 элементов")
    void testVarianceLessThanTwoElements() {
        List<Double> data = Collections.singletonList(5.0);
        assertThrows(RuntimeException.class, () -> loopAlgorithm.variance(data));
    }

    @Test
    @DisplayName("Мутация: skewness с очень малой дисперсией")
    void testSkewnessSmallVariance() {
        List<Double> data = Arrays.asList(1.0, 1.001, 1.002, 1.003);
        double skewness = loopAlgorithm.skewness(data);
        assertTrue(Double.isFinite(skewness));
    }

    @Test
    @DisplayName("Мутация: сравнение Loop vs Stream для всех методов")
    void testAllMethodsLoopVsStream() {
        List<Double> data = Arrays.asList(2.5, 3.5, 1.5, 4.5, 2.0);

        assertEquals(loopAlgorithm.min(data), streamAlgorithm.min(data), 0.0001);
        assertEquals(loopAlgorithm.max(data), streamAlgorithm.max(data), 0.0001);
        assertEquals(loopAlgorithm.mean(data), streamAlgorithm.mean(data), 0.0001);
        assertEquals(loopAlgorithm.median(data), streamAlgorithm.median(data), 0.0001);
        assertEquals(loopAlgorithm.variance(data), streamAlgorithm.variance(data), 0.0001);
        assertEquals(loopAlgorithm.skewness(data), streamAlgorithm.skewness(data), 0.0001);
        assertEquals(loopAlgorithm.kurtosis(data), streamAlgorithm.kurtosis(data), 0.0001);
    }

}

##StreamStatisticsAlgorithmTest
package algorithm;

import algorithms.LoopStatisticsAlgorithm;
import algorithms.StatisticsAlgorithm;
import algorithms.StreamStatisticsAlgorithm;
import org.junit.jupiter.api.Test;
import org.junit.jupiter.api.DisplayName;
import org.junit.jupiter.params.ParameterizedTest;
import org.junit.jupiter.params.provider.ValueSource;
import org.mockito.Mockito;

import static org.junit.jupiter.api.Assertions.*;
import static org.junit.jupiter.api.Assumptions.*;
import static org.hamcrest.MatcherAssert.assertThat;
import static org.hamcrest.Matchers.*;
import static org.mockito.Mockito.*;

import java.util.Arrays;
import java.util.Collections;
import java.util.List;

import static org.mockito.ArgumentMatchers.anyList;

@DisplayName("Тестирование StreamStatisticsAlgorithm")
public class StreamStatisticsAlgorithmTest {
    private final StatisticsAlgorithm loopAlgorithm = new LoopStatisticsAlgorithm();
    private final StatisticsAlgorithm streamAlgorithm = new StreamStatisticsAlgorithm();

    //ГРАНИЧНЫЕ ЗНАЧЕНИЯ (BoundaryValue)
    @Test
    @DisplayName("Граничное значение: пустой список должен выбрасывать исключение")
    void testMinEmptyList_BoundaryValue() {
        //граничное значение 0 элементов
        assertThrows(RuntimeException.class, () -> streamAlgorithm.min(Collections.emptyList()));
    }

    @Test
    @DisplayName("Граничное значение: один элемент")
    void testMinSingleElement_BoundaryValue() {
        //граничное значение 1 элемент
        List<Double> data = Collections.singletonList(5.0);
        double result = streamAlgorithm.min(data);
        assertEquals(5.0, result, 0.0001);
    }

    @Test
    @DisplayName("Граничное значение: 1 элемент для дисперсии должен выбрасывать исключение")
    void testVarianceOneElement_BoundaryValue() {
        //меньше минимального: 1 элемент
        List<Double> data = Collections.singletonList(5.0);
        assertThrows(RuntimeException.class, () -> streamAlgorithm.variance(data));
    }

    @Test
    @DisplayName("Граничное значение: 3 элемента для эксцесса должен выбрасывать исключение")
    void testKurtosisThreeElements_BoundaryValue() {
        //меньше минимального: 3 элемента
        List<Double> data = Arrays.asList(1.0, 2.0, 3.0);
        assertThrows(RuntimeException.class, () -> streamAlgorithm.kurtosis(data));
    }

    // КЛАССЫ ЭКВИВАЛЕНТНОСТИ (Equivalence Partitioning)
    @ParameterizedTest
    @ValueSource(doubles = {1.0, 5.5, 100.0})
    @DisplayName("Класс эквивалентности: положительные числа")
    void testMeanPositiveNumbers_EquivalenceClass(double value) {
        List<Double> data = Arrays.asList(value, value, value);
        double result = streamAlgorithm.mean(data);
        assertThat(result, is(closeTo(value, 0.0001)));
    }

    @ParameterizedTest
    @ValueSource(doubles = {-1.0, -5.5, -100.0})
    @DisplayName("Класс эквивалентности: отрицательные числа")
    void testMeanNegativeNumbers_EquivalenceClass(double value) {
        List<Double> data = Arrays.asList(value, value, value);
        double result = streamAlgorithm.mean(data);
        assertThat(result, is(closeTo(value, 0.0001)));
    }

    @Test
    @DisplayName("Медиана: нечетное количество (3 элемента)")
    void testMedianOddCount() {
        List<Double> data = Arrays.asList(1.0, 2.0, 3.0);
        assertEquals(2.0, streamAlgorithm.median(data), 0.0001);
    }

    @Test
    @DisplayName("Медиана: четное количество (4 элемента)")
    void testMedianEvenCount() {
        List<Double> data = Arrays.asList(1.0, 2.0, 3.0, 4.0);
        assertEquals(2.5, streamAlgorithm.median(data), 0.0001);
    }

    //ПОКРЫТИЕ ОПЕРАТОРОВ (Statement Testing)
    @Test
    @DisplayName("Покрытие операторов: все методы должны выполниться")
    void testAllMethodExecute_StatementCoverage() {
        List<Double> data = Arrays.asList(1.0, 2.0, 3.0, 4.0, 5.0);

        //каждый метод выполнится хотя бы один раз
        double min = streamAlgorithm.min(data);
        double max = streamAlgorithm.max(data);
        double mean = streamAlgorithm.mean(data);
        double median = streamAlgorithm.median(data);
        double variance = streamAlgorithm.variance(data);
        double skewness = streamAlgorithm.skewness(data);
        double kurtosis = streamAlgorithm.kurtosis(data);
        String name = streamAlgorithm.getName();

        //проверка, что все значения получены
        assertAll("Все методы выполнены успешно",
                () -> assertNotNull(min),
                () -> assertNotNull(max),
                () -> assertNotNull(mean),
                () -> assertNotNull(median),
                () -> assertNotNull(variance),
                () -> assertNotNull(skewness),
                () -> assertNotNull(kurtosis),
                () -> assertNotNull(name));
    }

    //ПОКРЫТИЕ ВЕТВЛЕНИЙ (Branch Testing)
    @Test
    @DisplayName("Покрытие ветвлений: медиана для нечетного количества (else)")
    void testMedianOddCount_BranchCoverage() {
        //n%2!=0
        List<Double> data = Arrays.asList(1.0, 3.0, 2.0);
        double result = streamAlgorithm.median(data);
        assertEquals(2.0, result, 0.0001);
    }

    @Test
    @DisplayName("Покрытие ветвлений: медиана для нечетного количества (if)")
    void testMedianEvenCount_BranchCoverage() {
        //n%2==0
        List<Double> data = Arrays.asList(1.0, 3.0, 2.0, 4.0);
        double result = streamAlgorithm.median(data);
        assertEquals(2.5, result, 0.0001);
    }

    @Test
    @DisplayName("Покрытие ветвлений: isEmpty() - true/false")
    void testIsEmpty_BranchCoverage() {
        //пустой список
        assertThrows(RuntimeException.class, () -> streamAlgorithm.min(Collections.emptyList()));

        //заполненный список
        List<Double> data = Collections.singletonList(1.0);
        assertDoesNotThrow(() -> streamAlgorithm.min(data));
    }

    //СРАВНЕНИЕ АЛГОРИТМОВ
    @Test
    @DisplayName("Сравнение алгоритмов: Loop и Stream должны давать одинаковые результаты")
    void testLoopVsStreamAlgorithmComparison() {
        List<Double> data = Arrays.asList(1.0, 5.0, 3.0, 9.0, 2.0, 7.0, 4.0);

        double loopMin = loopAlgorithm.min(data);
        double streamMin = streamAlgorithm.min(data);
        assertEquals(loopMin, streamMin, 0.0001, "min должен совпадать");

        double loopMax = loopAlgorithm.max(data);
        double streamMax = streamAlgorithm.max(data);
        assertEquals(loopMax, streamMax, 0.0001, "max должен совпадать");

    }

    //ASSUMPTIONS
    @Test
    @DisplayName("Assumption: тест имеет смысл только для положительных размеров")
    void testWithAssumptionOnSize() {
        List<Double> data = Arrays.asList(1.0, 2.0, 3.0);

        //предположение: тест имеет смысл только если данных больше 2
        assumeTrue(data.size() > 2, "Требуется минимум 3 элемента");

        double variance = streamAlgorithm.variance(data);
        assertTrue(variance >= 0, "Дисперсия не может быть отрицательной");
    }

    @Test
    @DisplayName("Assumption: пропускаем тест если выборка пустая")
    void testWithAssumptionOnEmptyList() {
        List<Double> data = Arrays.asList(1.0, 2.0, 3.0, 4.0, 5.0);

        //предположение: данные не должны быть пустыми
        assumeFalse(data.isEmpty(), "Данные не должны быть пустыми");

        double mean = streamAlgorithm.mean(data);
        assertThat(mean, is(greaterThan(0.0)));
    }

    //MOCKS
    @Test
    @DisplayName("Mock: создание мок объекта StatisticsAlgorithm")
    void testMockStatisticsAlgorithm() {
        StatisticsAlgorithm mockAlg = Mockito.mock(StatisticsAlgorithm.class);
        assertNotNull(mockAlg);

        double result = mockAlg.mean(Collections.emptyList());
        assertEquals(0.0, result, 0.0001);
    }

    @Test
    @DisplayName("Stub: настройка поведения мок объекта")
    void testStubStatisticsAlgorithm() {
        StatisticsAlgorithm stubAlg = Mockito.mock(StatisticsAlgorithm.class);

        when(stubAlg.mean(anyList())).thenReturn(42.0);
        when(stubAlg.min(anyList())).thenReturn(1.0);
        when(stubAlg.max(anyList())).thenReturn(100.0);

        assertEquals(42.0, stubAlg.mean(Collections.emptyList()), 0.0001);
        assertEquals(1.0, stubAlg.min(Collections.emptyList()), 0.0001);
        assertEquals(100.0, stubAlg.max(Collections.emptyList()), 0.0001);

        verify(stubAlg, times(1)).mean(anyList());
        verify(stubAlg, times(1)).min(anyList());
        verify(stubAlg, times(1)).max(anyList());
    }

    @Test
    @DisplayName("Spy: шпион на реальном объекте StreamStatisticsAlgorithm")
    void testSpyStatisticsAlgorithm() {
        StreamStatisticsAlgorithm spyAlg = Mockito.spy(new StreamStatisticsAlgorithm());

        List<Double> data = Arrays.asList(1.0, 2.0, 3.0, 4.0, 5.0);
        double mean = spyAlg.mean(data);

        verify(spyAlg, atLeastOnce()).mean(data);
        assertEquals(3.0, mean, 0.0001);
    }

    //MATCHERS
    @Test
    @DisplayName("Matchers: использование Hamcrest матчеров")
    void testWithHamcrestMatchers() {
        List<Double> data = Arrays.asList(1.0, 2.0, 3.0, 4.0, 5.0);

        double mean = streamAlgorithm.mean(data);
        double min = streamAlgorithm.min(data);
        double max = streamAlgorithm.max(data);

        //matcher 1 closeTo
        assertThat(mean, is(closeTo(3.0, 0.0001)));

        //matcher 2 greater/lessthan
        assertThat(min, lessThan(mean));
        assertThat(max, greaterThan(mean));

        //matcher 3 equalto
        assertThat(min, is(equalTo(1.0)));
        assertThat(max, is(equalTo(5.0)));
    }

    //НЕГАТИВНЫЕ СЦЕНАРИИ
    @Test
    @DisplayName("Негативный сценарий: null должен выбрасывать исключение")
    void testWithNullData() {
        assertThrows(NullPointerException.class, () -> streamAlgorithm.mean(null));
    }

    @Test
    @DisplayName("Негативный сценарий: эксцесс при нулевой дисперсии")
    void testKurtosisZeroVariance() {
        List<Double> data = Arrays.asList(5.0, 5.0, 5.0, 5.0); //дисперсия 0, т.к. элементы одинаковые
        assertThrows(RuntimeException.class, () -> streamAlgorithm.kurtosis(data));
    }
}

##JsonUtilsTest
package utils;

import json.SampleData;
import org.junit.jupiter.api.*;
import org.junit.jupiter.api.io.TempDir;
import org.junit.jupiter.params.ParameterizedTest;
import org.junit.jupiter.params.provider.ValueSource;
import org.mockito.internal.matchers.NotNull;

import static org.junit.jupiter.api.Assertions.*;
import static org.junit.jupiter.api.Assumptions.*;
import static org.hamcrest.MatcherAssert.assertThat;
import static org.hamcrest.Matchers.*;

import java.io.File;
import java.io.IOException;
import java.nio.file.Files;
import java.nio.file.Path;
import java.util.Arrays;
import java.util.List;

@DisplayName("Тестирование JsonUtils")
public class JsonUtilsTest {
    @TempDir
    Path tempDir; //временная директория для тестов

    //Boundary Value
    @Test
    @DisplayName("Граничное значение: пустая выборка")
    void testSaveAndLoadEmptyData_BoundaryValue() throws IOException {
        SampleData original = new SampleData("TestSample", Arrays.asList());
        Path testFile = tempDir.resolve("empty_test.json");

        //сохранение
        JsonUtils.saveToJson(original, testFile.toString());

        //проверка что файл создан
        assertTrue(Files.exists(testFile), "Файл должен быть создан");

        //загрузка
        SampleData loaded = JsonUtils.loadFromJson(testFile.toString());

        //проверка данных
        assertEquals(original.sampleType, loaded.sampleType);
        assertEquals(0, loaded.data.size(), "Размер данных должен быть больше нуля");
    }

    @Test
    @DisplayName("Граничное значение: один элемент")
    void testSaveAndLoadSingleElement_BoundaryValue() throws IOException {
        SampleData original = new SampleData("TestSample", Arrays.asList(42.0));
        Path testFile = tempDir.resolve("empty_test.json");


        JsonUtils.saveToJson(original, testFile.toString());
        SampleData loaded = JsonUtils.loadFromJson(testFile.toString());

        //проверка данных
        assertEquals(1, loaded.data.size());
        assertEquals(42.0, loaded.data.get(0), 0.0001);
    }

    @Test
    @DisplayName("Граничное значение: большой файл (1000 элементов)")
    void testSaveAndLoadLargeData() throws IOException{
        List<Double> largeData = Arrays.asList(1.0, 2.0, 3.0); //упрощенно
        SampleData original = new SampleData("TestSample", largeData);
        Path testFile = tempDir.resolve("large_test.json");

        JsonUtils.saveToJson(original, testFile.toString());
        SampleData loaded = JsonUtils.loadFromJson(testFile.toString());

        assertEquals(original.data.size(), loaded.data.size());
        assertThat(loaded.data.size(), greaterThan(0));
    }

    //EQUIVALENCE PARTITIONING
    @ParameterizedTest
    @ValueSource(strings = {"test1.json", "test2.json", "data.json"})
    @DisplayName("Класс эквивалентности: разные имена файлов")
    void testSaveAndLoadVariousFilenames_EquivalenceClass(String filename) throws IOException{
        SampleData original = new SampleData("UniformSample", Arrays.asList(1.0,2.0,3.0));
        Path testFile = tempDir.resolve(filename);

        JsonUtils.saveToJson(original, testFile.toString());
        SampleData loaded = JsonUtils.loadFromJson(testFile.toString());

        assertEquals(original.sampleType, loaded.sampleType);
        assertEquals(original.data.size(), loaded.data.size());
    }

    @ParameterizedTest
    @ValueSource(doubles = {-100.0, 100.0, 0.0, 3.14526})
    @DisplayName("Класс эквивалентности: разные значения в выборке")
    void testSaveAndLoadVariousValues_EquivalenceClass(double value) throws IOException{
        SampleData original = new SampleData("TestSample", Arrays.asList(value, value*2, value*3));
        Path testFile = tempDir.resolve("Values_test.json");

        JsonUtils.saveToJson(original, testFile.toString());
        SampleData loaded = JsonUtils.loadFromJson(testFile.toString());

        assertEquals(original.data.size(), loaded.data.size());
        for (int i = 0;i<original.data.size();i++) {
            assertEquals(original.data.get(i), loaded.data.get(i), 0.0001);
        }
    }


    //ПОКРЫТИЕ ОПЕРАТОРОВ STATEMENT TESTING
    @Test
    @DisplayName("Покрытие операторов: все методы JsonUtils должны выполниться")
    void testAllMethodExecute_StatementCoverage() throws IOException{
        SampleData testData = new SampleData("TestSample", Arrays.asList(1.0,2.0,3.0));
        Path testFile = tempDir.resolve("coverage_test.json");

        JsonUtils.saveToJson(testData, testFile.toString());
        SampleData loaded = JsonUtils.loadFromJson(testFile.toString());
        String jsonString = JsonUtils.toJson(testData);

        //проверка, что все значения получены
        assertAll("Все методы выполнены успешно",
                ()->assertNotNull(loaded),
                ()->assertNotNull(jsonString),
                ()->assertTrue(jsonString.length()>0,"JSON строка не должна быть пустой"),
                ()->assertThat(jsonString, containsString("TestSample"))
        );
    }

    //ПОКРЫТИЕ ВЕТВЛЕНИЙ BRANCH TESTING
    @Test
    @DisplayName("Покрытие ветвлений: файл существует после сохранения")
    void testFileExistsAfterSave_BranchCoverage() throws IOException{
        SampleData testData = new SampleData("TestSample", Arrays.asList(1.0,2.0));
        Path testFile = tempDir.resolve("exists_test.json");

        JsonUtils.saveToJson(testData, testFile.toString());

        assertTrue(Files.exists(testFile), "Файл должен существовать после сохранения");
        assertTrue(Files.isRegularFile(testFile), "Должен быть обычным файлом");
    }

    @Test
    @DisplayName("Покрытие ветвлений: загрузка несуществующего файла должна выбрасывать исключение")
    void testLoadNonExistentFile_BranchCoverage(){
        Path nonExistentFile = tempDir.resolve("non_existent.json");

        assertThrows(RuntimeException.class, ()->JsonUtils.loadFromJson(nonExistentFile.toString()));
    }

    //НЕГАТИВНЫЕ СЦЕНАРИИ
    @Test
    @DisplayName("Негативный сценарий: загрузка несуществующего файла")
    void testNonExistentFile_Negative(){
        assertThrows(RuntimeException.class, ()->JsonUtils.loadFromJson("non_existent_file.json"));
    }

    @Test
    @DisplayName("Негативный сценарий: сохранение в недопустимый путь")
    void testSaveToInvalidPath_Negative(){
        SampleData testData = new SampleData("TestSample", Arrays.asList(1.0,2.0));

        //попытка сохранения в несуществующую директорию
        assertThrows(RuntimeException.class,()->JsonUtils.saveToJson(testData, "/nonexistent/path/file.json"));

    }

    @Test
    @DisplayName("Негативный сценарий: загрузка невалидного Json")
    void testLoadInvalidJson_Negative() throws IOException{
        Path testFile = tempDir.resolve("invalid.json");
        Files.writeString(testFile, "not valid json {{{");
        assertThrows(RuntimeException.class, ()->JsonUtils.loadFromJson(testFile.toString()));
    }

    //ASSUMPTIONS
    @Test
    @DisplayName("Assumption: тест имеет смысл только если директория доступна для записи")
    void testWithAssumptionOnWritableDirectory() throws IOException{
        assumeTrue(Files.exists(tempDir), "Временная директория должна существовать");
        assumeTrue(Files.isWritable(tempDir), "Директория должна быть доступна для записи");

        SampleData testData = new SampleData("TestSample", Arrays.asList(1.0,2.0,3.0));
        Path testFile = tempDir.resolve("assumption_test.json");

        JsonUtils.saveToJson(testData, testFile.toString());
        assumeTrue(Files.exists(testFile));
    }

    @Test
    @DisplayName("Assumption: тест имеет смысл только для файлов с расширением json")
    void testWithAssumptionWithFileExtension() throws IOException{
        String filename = "test.json";

        assumeTrue(filename.endsWith(".json"), "Файл должен иметь расширение .json");

        SampleData testData = new SampleData("TestSample", Arrays.asList(1.0,2.0));
        Path testFile = tempDir.resolve(filename);

        JsonUtils.saveToJson(testData, testFile.toString());
        SampleData loaded = JsonUtils.loadFromJson(testFile.toString());

        assertEquals(testData.sampleType, loaded.sampleType);

    }

    //MATCHERS
    @Test
    @DisplayName("Matchers: использование Hamcrest матчеров")
    void testWithHamcrestMatchers() throws IOException{
        SampleData testData = new SampleData("UniformSample", Arrays.asList(1.0,2.0,3.0,4.0,5.0));
        Path testFile = tempDir.resolve("matchers_test.json");

        JsonUtils.saveToJson(testData, testFile.toString());
        String jsonString = JsonUtils.toJson(testData);

        assertThat(jsonString, containsString("UniformSample"));
        assertThat(jsonString, is(notNullValue()));
        assertThat(jsonString.length(), is(greaterThan(10)));
        SampleData loaded = JsonUtils.loadFromJson(testFile.toString());
        assertThat(loaded.sampleType, is(equalTo("UniformSample")));

    }

    //КРУГОВОЕ ТЕСТИРОВАНИЕ
    @Test
    @DisplayName("Круговое тестирование: сохранение-загрузка-сравнение")
    void testRoundTripSaveLoad() throws IOException{
        SampleData original = new SampleData("NormalSample", Arrays.asList(1.0,2.0,3.0,4.0,5.0)).
                withNormalParams(3.0,1.5);
        Path testFile = tempDir.resolve("roundtrip_test.json");

        JsonUtils.saveToJson(original, testFile.toString());
        SampleData loaded = JsonUtils.loadFromJson(testFile.toString());

        assertEquals(original.sampleType, loaded.sampleType, "Тип выборки должен совпадать");
        assertEquals(original.data.size(), loaded.data.size(), "Размер данных должен совпадать");

        for (int i = 0; i<original.data.size();i++){
            assertEquals(original.data.get(i), loaded.data.get(i), 0.0001,
                    "Элемент " + i + " должен совпадать");
        }

        assertEquals(original.meanVal, loaded.meanVal, 0.0001);
        assertEquals(original.stddevVal, loaded.stddevVal, 0.0001);
    }

    @Test
    @DisplayName("Круговое тестирование: toJson-парсинг-сравнение")
    void testToJsonAndParse(){
        SampleData original = new SampleData("TestSample", Arrays.asList(10.0, 20.0, 30.0));

        //json в строку
        String jsonString = JsonUtils.toJson(original);

        //проверка что строка валидная
        assertThat(jsonString, is(notNullValue()));
        assertThat(jsonString, containsString("TestSample"));
        assertThat(jsonString, containsString("10.0"));

    }

    @Test
    @DisplayName("Проверка: перезапись существующего файла")
    void testOverwriteExistingFile() throws IOException{
        Path testFile = tempDir.resolve("overwrite_test.json");

        //первое сохранение
        SampleData first = new SampleData("FirstSample", Arrays.asList(1.0, 2.0));
        JsonUtils.saveToJson(first, testFile.toString());

        //перезапись
        SampleData second = new SampleData("SecondSample", Arrays.asList(3.0,4.0,5.0));
        JsonUtils.saveToJson(second, testFile.toString());

        SampleData loaded = JsonUtils.loadFromJson(testFile.toString());

        assertEquals("SecondSample", loaded.sampleType);
        assertEquals(3, loaded.data.size());
    }


}

##StatsUtilsTest
package utils;

import algorithms.LoopStatisticsAlgorithm;
import algorithms.StreamStatisticsAlgorithm;
import org.junit.jupiter.api.Test;
import org.junit.jupiter.api.DisplayName;
import org.junit.jupiter.params.ParameterizedTest;
import org.junit.jupiter.params.provider.ValueSource;
import samples.*;

import static org.junit.jupiter.api.Assertions.*;
import static org.junit.jupiter.api.Assumptions.*;
import static org.hamcrest.MatcherAssert.assertThat;
import static org.hamcrest.Matchers.*;

import java.util.Arrays;
import java.util.Collections;
import java.util.List;

@DisplayName("Тестирование StatsUtils")
public class StatsUtilsTest {

    //ГРАНИЧНЫЕ ЗНАЧЕНИЯ (Boundary Value Analysis)
    @Test
    @DisplayName("Граничное значение: пустые выборки должны выбрасывать исключение")
    void testCorrelationEmptyData_BoundaryValue(){
        UniformSample sample1 = new UniformSample();
        UniformSample sample2 = new UniformSample();

        sample1.setData(Collections.emptyList());
        sample2.setData(Collections.emptyList());

        assertThrows(RuntimeException.class, ()->StatsUtils.correlation(sample1, sample2));
    }

    @Test
    @DisplayName("Граничное значение: один элемент в выборке должен выбрасывать исключение")
    void testCorrelationSingleElement_BoundaryValue(){
        UniformSample sample1 = new UniformSample();
        UniformSample sample2 = new UniformSample();

        sample1.setData(Arrays.asList(1.0));
        sample2.setData(Arrays.asList(2.0));

        assertThrows(RuntimeException.class, ()->StatsUtils.correlation(sample1, sample2));
    }

    @Test
    @DisplayName("Граничное значение: два элемента - минимум для корреляции")
    void testCorrelationTwoElements_BoundaryValue(){
        UniformSample sample1 = new UniformSample();
        UniformSample sample2 = new UniformSample();

        sample1.setData(Arrays.asList(1.0, 2.0));
        sample2.setData(Arrays.asList(1.0, 2.0));

        double result = StatsUtils.correlation(sample1, sample2);
        assertEquals(1.0, result, 0.0001, "Коэффициент корреляции должен быть 1.0 для идентичных выборок");
    }

    @Test
    @DisplayName("Граничное значение: разные размеры выборок должны выбрасывать исключение")
    void testCorrelationDifferentSizes_BoundaryValue(){
        UniformSample sample1 = new UniformSample();
        UniformSample sample2 = new UniformSample();

        sample1.setData(Arrays.asList(1.0, 2.0, 3.0));
        sample2.setData(Arrays.asList(1.0, 2.0));

        assertThrows(RuntimeException.class, ()->StatsUtils.correlation(sample1, sample2));
    }

    //КЛАССЫ ЭКВИВАЛЕНТНОСТИ (Equivalence Partitioning)
    @Test
    @DisplayName("Класс эквивалентности: положительная корреляция")
    void testCorrelationPositive_EquivalenceClass(){
        UniformSample sample1 = new UniformSample();
        UniformSample sample2 = new UniformSample();

        // Положительная корреляция: x растет, y растет
        sample1.setData(Arrays.asList(1.0, 2.0, 3.0, 4.0, 5.0));
        sample2.setData(Arrays.asList(2.0, 4.0, 6.0, 8.0, 10.0)); // y = 2x

        double result = StatsUtils.correlation(sample1, sample2);
        assertThat(result, is(closeTo(1.0, 0.0001)));
    }

    @Test
    @DisplayName("Класс эквивалентности: отрицательная корреляция")
    void testCorrelationNegative_EquivalenceClass(){
        UniformSample sample1 = new UniformSample();
        UniformSample sample2 = new UniformSample();

        // Отрицательная корреляция: x растет, y уменьшается
        sample1.setData(Arrays.asList(1.0, 2.0, 3.0, 4.0, 5.0));
        sample2.setData(Arrays.asList(10.0, 8.0, 6.0, 4.0, 2.0)); // y = 12 - 2x

        double result = StatsUtils.correlation(sample1, sample2);
        assertThat(result, is(closeTo(-1.0, 0.0001)));
    }

    @Test
    @DisplayName("Класс эквивалентности: нулевая дисперсия второй выборки")
    void testCorrelationZeroVariance_EquivalenceClass(){
        UniformSample sample1 = new UniformSample();
        UniformSample sample2 = new UniformSample();

        // Нулевая дисперсия: независимые переменные
        sample1.setData(Arrays.asList(1.0, 2.0, 3.0, 4.0, 5.0));
        sample2.setData(Arrays.asList(5.0, 5.0, 5.0, 5.0, 5.0)); // konstante

        assertThrows(RuntimeException.class, ()->StatsUtils.correlation(sample1, sample2));
    }

    @ParameterizedTest
    @ValueSource(ints = {2, 5, 10, 20})
    @DisplayName("Класс эквивалентности: разные размеры выборок для положительной корреляции")
    void testCorrelationVariousSizes_EquivalenceClass(int size){
        UniformSample sample1 = new UniformSample();
        UniformSample sample2 = new UniformSample();

        List<Double> data1 = new java.util.ArrayList<>();
        List<Double> data2 = new java.util.ArrayList<>();

        for (int i = 0; i < size; i++) {
            data1.add((double) i);
            data2.add((double) i * 2);
        }

        sample1.setData(data1);
        sample2.setData(data2);

        double result = StatsUtils.correlation(sample1, sample2);
        assertThat(result, is(closeTo(1.0, 0.0001)));
    }

    //ПОКРЫТИЕ ОПЕРАТОРОВ (Statement Testing)
    @Test
    @DisplayName("Покрытие операторов: все операции корреляции выполняются")
    void testAllStatements_StatementCoverage(){
        UniformSample sample1 = new UniformSample();
        UniformSample sample2 = new UniformSample();

        sample1.setData(Arrays.asList(1.0, 2.0, 3.0, 4.0, 5.0));
        sample2.setData(Arrays.asList(2.0, 4.0, 5.0, 8.0, 9.0));

        //выполнение корреляции охватит все операторы
        double correlation = StatsUtils.correlation(sample1, sample2);

        assertAll("Все операторы выполнены",
                ()->assertNotNull(correlation),
                ()->assertTrue(correlation >= -1.0 && correlation <= 1.0),
                ()->assertFalse(Double.isNaN(correlation))
        );
    }

    //ПОКРЫТИЕ ВЕТВЛЕНИЙ (Branch Testing)
    @Test
    @DisplayName("Покрытие ветвлений: проверка размера выборок (ветвь if)")
    void testBranchDifferentSizes_BranchCoverage(){
        UniformSample sample1 = new UniformSample();
        UniformSample sample2 = new UniformSample();

        sample1.setData(Arrays.asList(1.0, 2.0, 3.0));
        sample2.setData(Arrays.asList(1.0, 2.0));

        //ветвь: размеры не совпадают
        assertThrows(RuntimeException.class, ()->StatsUtils.correlation(sample1, sample2));
    }

    @Test
    @DisplayName("Покрытие ветвлений: проверка нулевой дисперсии первой выборки")
    void testBranchZeroVarianceFirst_BranchCoverage(){
        UniformSample sample1 = new UniformSample();
        UniformSample sample2 = new UniformSample();

        // Первая выборка имеет нулевую дисперсию (все значения одинаковые)
        sample1.setData(Arrays.asList(5.0, 5.0, 5.0, 5.0));
        sample2.setData(Arrays.asList(1.0, 2.0, 3.0, 4.0));

        assertThrows(RuntimeException.class, ()->StatsUtils.correlation(sample1, sample2));
    }

    @Test
    @DisplayName("Покрытие ветвлений: проверка нулевой дисперсии второй выборки")
    void testBranchZeroVarianceSecond_BranchCoverage(){
        UniformSample sample1 = new UniformSample();
        UniformSample sample2 = new UniformSample();

        // Вторая выборка имеет нулевую дисперсию
        sample1.setData(Arrays.asList(1.0, 2.0, 3.0, 4.0));
        sample2.setData(Arrays.asList(5.0, 5.0, 5.0, 5.0));

        assertThrows(RuntimeException.class, ()->StatsUtils.correlation(sample1, sample2));
    }

    //СРАВНЕНИЕ АЛГОРИТМОВ
    @Test
    @DisplayName("Сравнение алгоритмов: Loop vs Stream для корреляции")
    void testAlgorithmComparison(){
        // Создание выборок с Loop алгоритмом
        UniformSample sample1Loop = new UniformSample();
        UniformSample sample2Loop = new UniformSample();
        sample1Loop.setAlgorithm(new LoopStatisticsAlgorithm());
        sample2Loop.setAlgorithm(new LoopStatisticsAlgorithm());

        // Создание выборок со Stream алгоритмом
        UniformSample sample1Stream = new UniformSample();
        UniformSample sample2Stream = new UniformSample();
        sample1Stream.setAlgorithm(new StreamStatisticsAlgorithm());
        sample2Stream.setAlgorithm(new StreamStatisticsAlgorithm());

        List<Double> data1 = Arrays.asList(1.0, 2.0, 3.0, 4.0, 5.0);
        List<Double> data2 = Arrays.asList(2.0, 4.0, 6.0, 8.0, 10.0);

        sample1Loop.setData(data1);
        sample2Loop.setData(data2);
        sample1Stream.setData(data1);
        sample2Stream.setData(data2);

        double correlationLoop = StatsUtils.correlation(sample1Loop, sample2Loop);
        double correlationStream = StatsUtils.correlation(sample1Stream, sample2Stream);

        assertEquals(correlationLoop, correlationStream, 0.0001,
                "Результаты корреляции должны совпадать независимо от алгоритма");
    }

    //ASSUMPTIONS
    @Test
    @DisplayName("Assumption: тест имеет смысл только для положительных размеров")
    void testAssumptionPositiveSize(){
        UniformSample sample1 = new UniformSample();
        UniformSample sample2 = new UniformSample();

        sample1.setData(Arrays.asList(1.0, 2.0, 3.0, 4.0, 5.0));
        sample2.setData(Arrays.asList(2.0, 4.0, 6.0, 8.0, 10.0));

        //предположение: выборки имеют положительный размер
        assumeTrue(sample1.getData().size() > 0, "Выборка должна содержать данные");

        double result = StatsUtils.correlation(sample1, sample2);
        assertTrue(result >= -1.0 && result <= 1.0, "Корреляция должна быть в диапазоне [-1, 1]");
    }

    @Test
    @DisplayName("Assumption: тест имеет смысл только если выборки одного размера")
    void testAssumptionSameSizeRequired(){
        UniformSample sample1 = new UniformSample();
        UniformSample sample2 = new UniformSample();

        sample1.setData(Arrays.asList(1.0, 2.0, 3.0, 4.0));
        sample2.setData(Arrays.asList(5.0, 6.0, 7.0, 8.0));

        //предположение: размеры совпадают
        assumeTrue(sample1.getData().size() == sample2.getData().size(),
                "Выборки должны быть одного размера");

        double result = StatsUtils.correlation(sample1, sample2);
        assertNotNull(result);
    }

    //НЕГАТИВНЫЕ СЦЕНАРИИ
    @Test
    @DisplayName("Негативный сценарий: null в качестве первой выборки")
    void testNullFirstSample_NegativeScenario(){
        assertThrows(NullPointerException.class, ()->StatsUtils.correlation(null, new UniformSample()));
    }

    @Test
    @DisplayName("Негативный сценарий: null в качестве второй выборки")
    void testNullSecondSample_NegativeScenario(){
        assertThrows(NullPointerException.class, ()->StatsUtils.correlation(new UniformSample(), null));
    }

    @Test
    @DisplayName("Негативный сценарий: выборка с одинаковыми значениями в обеих")
    void testIdenticalValues_NegativeScenario(){
        UniformSample sample1 = new UniformSample();
        UniformSample sample2 = new UniformSample();

        // Обе выборки имеют нулевую дисперсию
        sample1.setData(Arrays.asList(5.0, 5.0, 5.0, 5.0));
        sample2.setData(Arrays.asList(3.0, 3.0, 3.0, 3.0));

        assertThrows(RuntimeException.class, ()->StatsUtils.correlation(sample1, sample2));
    }

    //MATCHERS
    @Test
    @DisplayName("Matchers: проверка диапазона корреляции")
    void testCorrelationRange_Matchers(){
        UniformSample sample1 = new UniformSample();
        UniformSample sample2 = new UniformSample();

        sample1.setData(Arrays.asList(1.0, 2.0, 3.0, 4.0, 5.0));
        sample2.setData(Arrays.asList(2.0, 3.0, 5.0, 8.0, 9.0));

        double result = StatsUtils.correlation(sample1, sample2);

        //matcher: результат в диапазоне [-1, 1]
        assertThat(result, is(greaterThanOrEqualTo(-1.0)));
        assertThat(result, is(lessThanOrEqualTo(1.0)));
    }

    @Test
    @DisplayName("Matchers: проверка идентичных выборок")
    void testIdenticalSamples_Matchers(){
        UniformSample sample1 = new UniformSample();
        UniformSample sample2 = new UniformSample();

        sample1.setData(Arrays.asList(1.0, 2.0, 3.0, 4.0, 5.0));
        sample2.setData(Arrays.asList(1.0, 2.0, 3.0, 4.0, 5.0));

        double result = StatsUtils.correlation(sample1, sample2);

        //matcher: идеальная положительная корреляция
        assertThat(result, is(closeTo(1.0, 0.0001)));
    }

    //ДОПОЛНИТЕЛЬНЫЕ ТЕСТЫ
    @Test
    @DisplayName("Сценарий: слабая положительная корреляция")
    void testWeakPositiveCorrelation(){
        NormalSample sample1 = new NormalSample(0, 1);
        NormalSample sample2 = new NormalSample(0, 1);

        sample1.setData(Arrays.asList(1.0, 2.0, 3.0, 4.0, 5.0));
        sample2.setData(Arrays.asList(1.5, 2.2, 3.1, 3.9, 5.1));

        double result = StatsUtils.correlation(sample1, sample2);

        assertThat(result, is(greaterThan(0.0)));
        assertThat(result, is(lessThan(1.0)));
    }

    @Test
    @DisplayName("Сценарий: умеренная положительная корреляция")
    void testModeratePositiveCorrelation(){
        UniformSample sample1 = new UniformSample();
        UniformSample sample2 = new UniformSample();

        sample1.setData(Arrays.asList(1.0, 2.0, 3.0, 4.0, 5.0, 6.0));
        sample2.setData(Arrays.asList(1.0, 3.0, 2.0, 4.0, 5.0, 6.0));

        double result = StatsUtils.correlation(sample1, sample2);

        assertTrue(result > 0.5 && result < 1.0, "Корреляция должна быть умеренной");
    }
}

\end{verbatim}
